\documentclass[11pt]{article}

\usepackage[margin=1in]{geometry}
\usepackage[T1]{fontenc}
\usepackage{lmodern}
\usepackage{amsmath,amssymb,amsthm,mathtools}
\usepackage{booktabs}
\usepackage{array}
\usepackage{microtype}
\usepackage[hidelinks]{hyperref}

\theoremstyle{plain}
\newtheorem{theorem}{Theorem}[section]
\newtheorem{proposition}[theorem]{Proposition}
\newtheorem{lemma}[theorem]{Lemma}
\newtheorem{corollary}[theorem]{Corollary}

\theoremstyle{definition}
\newtheorem{definition}[theorem]{Definition}
\theoremstyle{remark}
\newtheorem{remark}[theorem]{Remark}

\newcommand{\R}{\mathbb{R}}
\newcommand{\dd}{\mathrm{d}}
\newcommand{\eps}{\varepsilon}
\newcommand{\Ibar}{\bar{I}}
\newcommand{\wbar}{\bar{\omega}}

\title{\textbf{A Machine-Checked Equation-of-State-Independent Kerr Lower Bound on the Moment of Inertia of Relativistic Stars}}
\author{Serghei Iancenco\thanks{Independent researcher. Correspondence via \url{https://github.com/iancenco/relativistic-bounds-lean4}. Methodological disclosure: the continuum certificate search, Lean~4 formalization, and statement-level verification audit were conducted with an AI-assisted multi-model workflow directed by the author; see Section~\ref{sec:lean}.}}
\date{November 23, 2026}

\begin{document}
\maketitle

\begin{abstract}
We prove an equation-of-state-independent lower bound on the slow-rotation moment of inertia of static, spherically symmetric, horizonless isotropic relativistic stars in four-dimensional general relativity. Specifically, for any piecewise-smooth star of radius $R>0$ and gravitational mass $M>0$ obeying the Tolman--Oppenheimer--Volkoff (TOV) equations and Hartle's first-order slow-rotation equation off a finite set of interior density discontinuities, with non-increasing mean energy density ($4\pi r^3 \eps(r) \le 3 m(r)$, implied in particular by non-increasing local density $\eps$), isotropic pressure satisfying the dominant energy condition ($0 \le p \le \eps$), and a vacuum surface $p(R)=0$, the exterior Hartle moment of inertia strictly exceeds the Kerr black-hole value:
\[
  I \;>\; 4 M^3 \;=\; I_{\mathrm{Kerr}} \qquad (G = c = 1).
\]
No nuclear equation of state and no subluminal sound-speed assumption are required. The proof combines an auxiliary Bondi--Pontryagin compactness lemma ($m(r)/r < 51/125 = 0.408$) with an exact shifted Hartle--TOV Riccati identity along the Schwarzschild vacuum curve, discharged entirely by exact rational Bernstein positivity certificates and barrier propagation across finitely many interfaces. The complete analytical proof chain from the reduced TOV--Hartle system to $I > 4 M^3$ is machine-checked in Lean~4.36.0 + Mathlib with zero \texttt{sorry}, zero \texttt{native\_decide}, and only the three standard Lean axioms. We also provide explicit analytic non-vacuity witnesses (smooth and two-layer discontinuous), place the theorem alongside prior bounds (Hartle, Bondi, Barraco--Hamity, Karageorgis--Stalker, Andr\'easson, and Barta--Gergely), clarify why $4 M^3$ is the Kerr benchmark rather than the sharp infimum of the class, and show why the isotropy hypothesis is essential.
\end{abstract}

\tableofcontents

\section{Introduction and Literature Context}
\label{sec:intro}

In geometrized units ($G = c = 1$), the moment of inertia of a slowly rotating Kerr black hole of mass $M$, defined by $I_{\mathrm{Kerr}} \equiv \lim_{a \to 0} J/\Omega_H$ where $J = M a$ and $\Omega_H = a/(2 M r_+)$, takes the universal value
\begin{equation}
  I_{\mathrm{Kerr}} \;=\; 4 M^3, \qquad \Ibar_{\mathrm{Kerr}} \;\equiv\; \frac{I_{\mathrm{Kerr}}}{M^3} \;=\; 4.
  \label{eq:kerr-inertia}
\end{equation}
Numerical studies of relativistic stellar models have long observed that the dimensionless moment of inertia $\Ibar = I/M^3$ of isotropic fluid stars lies strictly above the Kerr value $4$, approaching $\Ibar \approx 4.046$ only in the extreme incompressible Buchdahl limit $M/R \to 4/9$ \cite{ChandrasekharMiller1974,Posada2017} and remaining well above $4$ across realistic nuclear equations of state \cite{YagiYunes2013,BreuRezzolla2016}. Recently, Barta and Gergely \cite{BartaGergely2026} explicitly highlighted $\Ibar > 4$ across four tabulated nuclear equations of state in their Hartle--Thorne analysis of neutron-star spin flip-flops.

Despite this strong numerical consensus, rigorous equation-of-state-independent lower bounds on $I/M^3$ have received far less analytical attention than compactness bounds. The classical work of Sabbadini and Hartle \cite{SabbadiniHartle1977} and Hartle \cite{Hartle1978} (especially Section~5.3.2 of \cite{Hartle1978}) established variational upper and lower bounds on the mass and moment of inertia of relativistic stars under the assumption that the equation of state is known up to a fiducial matching density $\rho_0$ and satisfies causality ($\dd p/\dd \eps \le 1$) above $\rho_0$. Those bounds are tight and astrophysically informative for neutron stars, but they depend on the low-density equation of state and do not establish a universal global inequality $\Ibar > 4$ independent of all microphysics.

The purpose of this paper is to give a short, self-contained, and completely machine-checked proof of the strict Kerr lower bound $I > 4 M^3$ for piecewise-smooth, horizonless, isotropic perfect-fluid stars satisfying the dominant energy condition ($0 \le p \le \eps$) and non-increasing mean density ($4\pi r^3 \eps \le 3 m$, which is implied by non-increasing local density $\eps$). No equation of state is specified anywhere in the star, no sound-speed causality constraint is assumed, and finitely many first-order interior density discontinuities (phase transitions) as well as a finite surface density jump are permitted.

\paragraph{Relation to compactness theorems.}
Our proof requires an upper bound on the local compactness $u(r) = m(r)/r$ so that the metric factor $s(r)^2 = 1 - 2 m(r)/r$ stays strictly above $23/125$. For this purpose we prove a self-contained auxiliary lemma, $m(r)/r < 51/125 = 0.408$, using a two-piece polynomial barrier in Bondi's phase plane \cite{Bondi1964}. We emphasize that this compactness bound is an \emph{auxiliary ingredient}, not a new sharp compactness record: Buchdahl \cite{Buchdahl1959} proved $M/R \le 4/9$ for non-increasing isotropic fluids without energy conditions, Andr\'easson \cite{Andreasson2008} proved $2M/R \le 8/9$ for anisotropic spheres under $p_r + 2p_t \le \eps$, and Bondi \cite{Bondi1964} and Karageorgis--Stalker \cite{KarageorgisStalker2008} analyzed sharp compactness ceilings under $p \le w\eps$.

It is worth clarifying a subtlety regarding the bound $M/R \le 3/8 = 0.375$ often associated with uniform-density stars and Barraco--Hamity \cite{BarracoHamity2002}. For a \emph{constant-density} isotropic star ($\eps = \bar{\eps}$), Buchdahl's central-pressure formula gives $p_c/\bar{\eps} = (1-\sqrt{1-2C})/(3\sqrt{1-2C}-1)$, so the dominant energy condition $p_c \le \bar{\eps}$ indeed forces $C = M/R \le 3/8 = 0.375$. For a \emph{non-uniform} monotone-density star, however, Buchdahl's comparison bounds $p_c/\bar{\eps}$ from below while DEC only constrains $p_c \le \eps_c$, and $\eps_c > \bar{\eps}$ allows $p_c/\bar{\eps} > 1$. Indeed, as verified in Section~\ref{sec:sharpness}, an explicit two-zone monotone isotropic DEC star consisting of a stiff $p = \eps$ core continuously matched to a constant-density envelope attains $M/R = 0.398800 > 3/8$ ($2M/R = 0.797600$), in agreement with Bondi's analysis \cite{Bondi1964}. Thus $3/8$ does not bound the full monotone isotropic DEC class, whereas our auxiliary polynomial barrier $m/r < 51/125 = 0.408$ sits within $2.3\%$ of this explicit realization and yields a short, zero-floating-point proof in Lean~4.

\paragraph{Scope, sharpness, and formalization honesty.}
We highlight three caveats up front:
\begin{enumerate}
  \item \emph{Reduced-ODE input to Lean.} The Lean~4 development formalizes the theorem starting from the reduced TOV equations, Hartle's Riccati equation for $Z = r s \wbar'/\wbar$, interface continuity conditions, and the matched surface inertia formula $I = R^3 Z(R)/(6 s(R) + 2 Z(R))$. The derivation of these reduced equations from the four-dimensional Einstein equations is standard and is given in Section~\ref{sec:setup}, but is not formalized from scratch in Lean.
  \item \emph{Non-vacuity.} Explicit smooth and two-layer discontinuous physical stars satisfying every Lean hypothesis are constructed and proved analytically in Section~\ref{sec:nonvacuity}, but formalizing an existence proof for Hartle's ODE inside Lean~4 + Mathlib remains future work.
  \item \emph{Sharpness and hypotheses.} The constant $4$ is the Kerr benchmark rather than the sharp infimum of the monotone DEC isotropic class (Section~\ref{sec:sharpness}): uniform DEC stars reach $\Ibar \approx 4.6549$ at $C=3/8$ and two-zone monotone DEC stars reach $\Ibar \approx 4.2855$. Relaxing isotropy allows strictly-DEC stars with $\Ibar < 4$, showing that isotropy is essential. Dropping DEC entirely while keeping monotone isotropic fluids (the full Buchdahl class, where uniform stars approach $\Ibar \approx 4.046$) remains an open problem.
\end{enumerate}

\section{Physical Setup, Interface Matching, and Bridge to Formalization}
\label{sec:setup}

\subsection{Derivation of the Reduced TOV--Hartle System}

In geometrized units ($G = c = 1$), consider a static, spherically symmetric, horizonless spacetime with metric
\begin{equation}
  \dd s^2 \;=\; -N(r)^2\,\dd t^2 + s(r)^{-2}\,\dd r^2 + r^2\left(\dd\theta^2 + \sin^2\theta\,\dd\phi^2\right),
  \qquad s(r) \;\equiv\; \sqrt{1 - \frac{2 m(r)}{r}} \;>\; 0,
  \label{eq:metric}
\end{equation}
populated by an isotropic perfect fluid with energy density $\eps(r)$ and pressure $p(r)$ supported on $[0, R]$, with vacuum exterior ($r > R$) and surface boundary condition $p(R) = 0$. The mixed $tt$ and $rr$ Einstein equations together with radial stress-energy conservation $\nabla_\mu T^{\mu}{}_{r} = 0$ give the standard TOV system \cite{Hartle1967}:
\begin{equation}
  m' \;=\; 4\pi r^2 \eps, \qquad
  \frac{N'}{N} \;=\; \frac{m + 4\pi r^3 p}{r^2 s^2}, \qquad
  p' \;=\; -(\eps + p)\frac{m + 4\pi r^3 p}{r^2 s^2}.
  \label{eq:tov}
\end{equation}
Introduce the dimensionless Bondi variables \cite{Bondi1964}
\begin{equation}
  u(r) \;\equiv\; \frac{m(r)}{r}, \qquad
  \pi_e(r) \;\equiv\; 4\pi r^2 \eps(r), \qquad
  \pi_p(r) \;\equiv\; 4\pi r^2 p(r),
  \label{eq:bondi-vars}
\end{equation}
so that $s^2 = 1 - 2u$ and \eqref{eq:tov} takes the logarithmic-radius form
\begin{equation}
  r u' \;=\; \pi_e - u, \qquad
  r s' \;=\; \frac{u - \pi_e}{s}, \qquad
  r \pi_p' \;=\; 2 \pi_p - \frac{(\pi_e + \pi_p)(u + \pi_p)}{s^2}.
  \label{eq:tov-bondi}
\end{equation}

Under rigid slow rotation at uniform angular velocity $\Omega$ to first order in $\Omega$, the metric acquires a single off-diagonal frame-dragging component $g_{t\phi} = -\omega(r) r^2 \sin^2\theta$ \cite{Hartle1967}. Defining the relative angular velocity $\wbar(r) \equiv \Omega - \omega(r)$ and the auxiliary function $j(r) \equiv s(r)/N(r)$, the $t\phi$ Einstein equation is the linear second-order ODE
\begin{equation}
  \frac{\dd}{\dd r}\!\left(r^4 j\,\wbar'\right) + 4 r^3 j'\,\wbar \;=\; 0,
  \label{eq:hartle-2nd}
\end{equation}
where from \eqref{eq:tov} the logarithmic derivative of $j$ obeys
\begin{equation}
  \frac{r j'}{j} \;=\; \frac{r s'}{s} - \frac{r N'}{N} \;=\; \frac{u - \pi_e}{s^2} - \frac{u + \pi_p}{s^2} \;=\; -\frac{\pi_e + \pi_p}{s^2}.
  \label{eq:j-deriv}
\end{equation}
Where $\wbar(r) > 0$, define $y(r) \equiv r \wbar'(r)/\wbar(r)$ and the Hartle Riccati variable
\begin{equation}
  Z(r) \;\equiv\; s(r)\,y(r) \;=\; \frac{r s(r)\,\wbar'(r)}{\wbar(r)}.
  \label{eq:Z-def}
\end{equation}
Expanding \eqref{eq:hartle-2nd} yields $r y' = -y^2 - 3y - (y+4)(r j'/j) = -y^2 - 3y + (\pi_e + \pi_p)(y+4)/s^2$. Differentiating $Z = s y$ via \eqref{eq:tov-bondi} gives Hartle's slow-rotation Riccati equation:
\begin{equation}
  r Z' \;=\; \frac{(u + \pi_p) Z}{s^2} - \frac{Z^2}{s} - 3 Z + \frac{4(\pi_e + \pi_p)}{s}.
  \label{eq:hartle-Z}
\end{equation}

\subsection{Exterior Matching and Interface Junction Conditions}

In the exterior vacuum ($r > R$), $j(r) = 1$ (normalizing $N(r) = s(r)$ at infinity) and \eqref{eq:hartle-2nd} integrates immediately to
\begin{equation}
  \wbar(r) \;=\; \Omega - \frac{2 J}{r^3}, \qquad J \;=\; I \Omega,
  \label{eq:vacuum-wbar}
\end{equation}
with angular momentum $J$ and moment of inertia $I$. Continuity of $\wbar$ and $\wbar'$ at $r = R$ (since $j$ is continuous across the vacuum surface) gives
\begin{equation}
  y(R) \;=\; \frac{R \wbar'(R)}{\wbar(R)} \;=\; \frac{6 I}{R^3 - 2 I}, \qquad
  Z(R) \;=\; s(R)\,y(R) \;=\; \frac{6 s(R)\,I}{R^3 - 2 I}.
  \label{eq:Z-surface}
\end{equation}
Solving \eqref{eq:Z-surface} for $I$ yields the exterior moment-of-inertia formula
\begin{equation}
  I \;=\; \frac{R^3 Z(R)}{6 s(R) + 2 Z(R)}.
  \label{eq:I-formula}
\end{equation}
Setting $I = \Ibar_0 M^3$ with $M = m(R)$ and $u_R = M/R = (1 - s(R)^2)/2$ in \eqref{eq:Z-surface} defines the \emph{Schwarzschild vacuum curve} at fixed dimensionless inertia $\Ibar_0$:
\begin{equation}
  Z_c(\Ibar_0, s) \;\equiv\; \frac{3 \Ibar_0\,s\,(1 - s^2)^3}{4 - \Ibar_0\,(1 - s^2)^3}.
  \label{eq:Zc-def}
\end{equation}
At $\Ibar_0 = 4$ (the Kerr value), $4 - 4(1 - s^2)^3 = 4(1 - 8 u^3) > 0$ for every $s \in (0, 1]$, and whenever $Z(R) > Z_c(4, s(R)) \ge 0$ we immediately have $6 s(R) + 2 Z(R) > 0$ and $I > 4 M^3$.

\paragraph{Interior density discontinuities (first-order phase transitions).}
Suppose the star has finitely many interface radii $r_k \in \mathcal{J} \subset (0, R]$ where $\eps$ jumps downward (including optionally the surface $R$ when $\eps(R^-) > 0$). By the Israel--Darmois junction conditions for an ordinary perfect-fluid interface with no singular surface layer, $m$, $p$, $N$, $s$, $\wbar$, and the angular-momentum flux $r^4 j \wbar'$ are all continuous across every $r_k \in \mathcal{J}$. Consequently,
\[
  Z \;=\; \frac{s \cdot (r^4 j \wbar')}{r^3 j \wbar}
\]
is continuous on $(0, R]$ even though $m'$, $p'$, and $Z'$ suffer finite jumps at $r_k \in \mathcal{J}$. On $(0, R] \setminus \mathcal{J}$, the ODEs \eqref{eq:tov} and \eqref{eq:hartle-Z} hold pointwise.

\subsection{Formalized Star Class and Mapping to Lean~4}

Table~\ref{tab:bridge} records the exact correspondence between the physical assumptions above and the Lean~4 structures \texttt{PiecewiseRegularIsotropicDECStar} and \texttt{MonotoneIsotropicDECStar} in \texttt{NaturalMechanics/T4A18\_KerrInertia.lean}.

\begin{table}[ht]
\centering
\small
\begin{tabular}{@{}p{0.34\textwidth}p{0.28\textwidth}p{0.32\textwidth}@{}}
\toprule
Physical assumption & Mathematical expression & Lean~4 field / hypothesis \\
\midrule
Positive radius and mass & $R > 0$, $m(0)=0$, $M=m(R)>0$ & \texttt{R\_pos}, \texttt{m\_zero}, \texttt{M\_pos} \\
Horizonless interior & $s(r) = \sqrt{1 - 2 m(r)/r} > 0$ on $(0, R]$ & \texttt{s\_pos} \\
Finite interface set & $\mathcal{J} \subset \R$ finite & \texttt{jumps : Finset $\R$} \\
Israel--Darmois matching & $m, p, Z$ continuous on $(0, R]$ & \texttt{continuousOn\_m}, \texttt{\_p}, \texttt{\_Z} \\
TOV \& Hartle ODEs off $\mathcal{J}$ & Eqs.~\eqref{eq:tov} and \eqref{eq:hartle-Z} on $(0,R]\setminus\mathcal{J}$ & \texttt{hasDerivAt\_m}, \texttt{\_p}, \texttt{\_Z} \\
Vacuum surface & $p(R) = 0$ & \texttt{p\_surface} \\
Weighted central limits & $m/r \to 0$, $4\pi r^2 p \to 0$, $Z \to 0$ as $r \to 0^+$ & \texttt{center\_u}, \texttt{center\_p}, \texttt{center\_Z}, \texttt{Z\_zero} \\
Isotropic DEC & $0 \le p(r) \le \eps(r)$ on $(0, R]$ & \texttt{p\_nonneg}, \texttt{p\_le\_$\eps$} \\
Non-increasing mean density & $4\pi r^3 \eps(r) \le 3 m(r)$ on $(0, R]$ & \texttt{mean\_density} \\
Non-increasing local density (stronger) & $\eps$ antitone on $(0, R]$ & \texttt{AntitoneOn S.$\eps$ (Ioc 0 S.R)} \\
\bottomrule
\end{tabular}
\caption{Exact correspondence between the physical model and the Lean~4 star structures. Note that the reduced ODEs and weighted central limits are taken as inputs to the formal proof; full $C^2$ regularity at $r=0$ is strictly stronger than the weighted central limits.}
\label{tab:bridge}
\end{table}

\begin{lemma}[One-way monotonicity implication across jumps]
\label{lem:mean-density}
Let $m, \eps : \R \to \R$ satisfy $m(0) = 0$, $\lim_{r\to 0^+} m(r)/r = 0$, $m$ continuous on $(0, R]$, and $m'(r) = 4\pi r^2 \eps(r)$ for all $r \in (0, R] \setminus \mathcal{J}$ with $\mathcal{J}$ finite. If $\eps$ is non-increasing on $(0, R]$, then
\begin{equation}
  4\pi r^3 \eps(r) \;\le\; 3 m(r) \qquad \forall\, r \in (0, R],
  \label{eq:mean-density}
\end{equation}
equivalently $\pi_e(r) \le 3 u(r)$ and $m(r)/r^3$ is non-increasing on $(0, R]$. The converse is false in general: a non-increasing mean density does not imply that local density $\eps$ is non-increasing.
\end{lemma}

\begin{proof}
Fix $r \in (0, R]$ and set $g(\rho) \equiv m(\rho) - \frac{4\pi}{3}\rho^3 \eps(r)$ on $[0, r]$. Since $m(\rho) = \rho (m(\rho)/\rho) \to 0 = m(0)$ as $\rho \to 0^+$, $g$ is continuous on $[0, r]$. For $\rho \in (0, r) \setminus \mathcal{J}$, $g'(\rho) = 4\pi \rho^2 (\eps(\rho) - \eps(r)) \ge 0$ because $\eps$ is non-increasing. By induction on the finite exceptional set $\mathcal{J} \cap (0, r)$, $g$ is non-decreasing on $[0, r]$, hence $g(r) \ge g(0) = 0$, giving \eqref{eq:mean-density}. For the failure of the converse, consider any smooth positive perturbation adding a small local density bump near an outer radius while keeping $\eps(r)$ strictly below the interior average $3 m(r)/(4\pi r^3)$.
\end{proof}

\section{Self-Contained Continuum Proof}
\label{sec:proof}

\subsection{Part A: Exact Shifted Hartle--TOV Riccati Identity}

For any real shift parameter $\alpha \in \R$ and reference inertia $\Ibar_0 \in \R$, define the shifted Hartle variable and its residual relative to the vacuum curve \eqref{eq:Zc-def}:
\begin{equation}
  W_\alpha(r) \;\equiv\; Z(r) - \frac{2\alpha\,\pi_p(r)}{s(r)}, \qquad
  D_{\alpha,\Ibar_0}(r) \;\equiv\; W_\alpha(r) - Z_c(\Ibar_0, s(r)).
  \label{eq:W-D-def}
\end{equation}

\begin{proposition}[Exact shifted Hartle--TOV Riccati identity]
\label{prop:partA}
Off the interface set $\mathcal{J}$, wherever $s > 0$ and $4 - \Ibar_0(1 - s^2)^3 \neq 0$, the residual $D = W_\alpha - Z_c(\Ibar_0, s)$ obeys the exact identity
\begin{equation}
  r D' \;=\; L\,D \;+\; c_e(\alpha, \Ibar_0, s)\,\pi_e \;+\; c_p(\alpha, \Ibar_0, s)\,\pi_p \;+\; c_{pp}(\alpha, s)\,\pi_p^2,
  \label{eq:partA}
\end{equation}
where the linear coefficient and matter coefficients are
\begin{align}
  L &\;=\; -\frac{W_\alpha + Z_c(\Ibar_0, s)}{s} + \frac{u + (1 - 4\alpha)\pi_p}{s^2} - 3, \label{eq:L-def} \\
  c_e(\alpha, \Ibar_0, s) &\;=\; \frac{N_e(\alpha, \Ibar_0, s^2)}{s^3\bigl(4 - \Ibar_0(1 - s^2)^3\bigr)^2}, \qquad
  c_p(\alpha, \Ibar_0, s) \;=\; \frac{N_p(\alpha, \Ibar_0, s^2)}{s^3\bigl(4 - \Ibar_0(1 - s^2)^3\bigr)}, \label{eq:ce-cp-def} \\
  c_{pp}(\alpha, s) &\;=\; \frac{4\alpha(1 - \alpha)}{s^3}, \label{eq:cpp-def}
\end{align}
with explicit polynomials in $t = s^2$:
\begin{align}
  N_e(\alpha, \Ibar_0, t) &\;=\; -\Bigl[
    \Ibar_0^2\bigl(\alpha(t-1)^7 - t(t-1)^6\bigr)
    + 8\Ibar_0\alpha(t-1)^4 \nonumber\\
  &\qquad + 4\Ibar_0\bigl(13 t^4 - 21 t^3 + 3 t^2 + 5 t\bigr)
    + 16\alpha(t-1) - 64 t
  \Bigr], \label{eq:NumE}\\
  N_p(\alpha, \Ibar_0, t) &\;=\; -\Bigl[
    \Ibar_0\alpha(t-1)^3(t-3) - \Ibar_0 t(t-1)^3 + 52\alpha t - 12\alpha - 16 t
  \Bigr]. \label{eq:NumP}
\end{align}
In particular, along the vacuum curve $D = 0$, both the matter-free constant term and the $\pi_e \pi_p$ cross term vanish identically for all $(\alpha, \Ibar_0)$.
\end{proposition}

\begin{proof}
Substitute $Z = W_\alpha + 2\alpha \pi_p/s$ and $u = (1 - s^2)/2$ into
\[
  r D' \;=\; r Z' - 2\alpha\left(\frac{r \pi_p'}{s} - \frac{\pi_p\,r s'}{s^2}\right) - \frac{\partial Z_c}{\partial s}(\Ibar_0, s)\,r s',
\]
using \eqref{eq:tov-bondi} and \eqref{eq:hartle-Z}. Grouping the terms proportional to $D = W_\alpha - Z_c$ via $W_\alpha^2 - Z_c^2 = (W_\alpha + Z_c) D$ isolates $L\,D$. The vacuum curve $Z_c$ is by construction an exact solution of the vacuum ($\pi_e = \pi_p = 0$) Hartle equation, so the matter-free term cancels identically. The coefficient of $\pi_e \pi_p$ from $-2\alpha(r\pi_p'/s - \pi_p r s'/s^2)$ is $2\alpha(1/s^3 - 1/s^3) = 0$, so the $\pi_e \pi_p$ cross term cancels identically as well. Collecting the remaining terms over the common denominator $s^3(4 - \Ibar_0(1-s^2)^3)^2$ yields \eqref{eq:ce-cp-def}--\eqref{eq:NumP} by direct polynomial algebra (verified symbolically in Lean by \texttt{field\_simp; ring} in \texttt{partA}).
\end{proof}

\subsection{Part B: Auxiliary Bondi--Pontryagin Compactness Lemma}

\begin{proposition}[Auxiliary compactness lemma]
\label{prop:partB}
Every star in \texttt{MonotoneIsotropicDECStar} satisfies
\begin{equation}
  u(r) \;=\; \frac{m(r)}{r} \;<\; \frac{51}{125} \;=\; 0.408 \qquad \text{and} \qquad
  s(r)^2 \;=\; 1 - 2 u(r) \;>\; \frac{23}{125} \;=\; 0.184
  \label{eq:compactness-wall}
\end{equation}
for all $r \in (0, R]$.
\end{proposition}

\begin{proof}
First, since $4\pi r^3 \eps \le 3 m$, the quotient $m(r)/r^3$ has non-positive derivative off $\mathcal{J}$ and is continuous on $(0, R]$, hence non-increasing on $(0, R]$. Because $m(R) > 0$, we obtain $m(r) > 0$ and $u(r) > 0$ on $(0, R]$.

In Bondi's $(u, v)$ plane with $v \equiv \pi_p = 4\pi r^2 p$, the conditions $0 \le p \le \eps$ and $4\pi r^3 \eps \le 3 m$ restrict the density control $e \equiv \pi_e$ to the compact interval $e \in [v, 3u]$ at every $r \in (0, R]$. Define the continuous two-piece polynomial barrier $B(u)$ joined at $u_j = 49/200$:
\begin{align}
  B_1(u) &\;=\; u\!\left(\frac{61}{20} - \frac{103}{20} u + \frac{4}{5} u^2 - \frac{103}{4} u^3\right),
  \qquad u \in \left[0, \frac{49}{200}\right], \label{eq:B1}\\
  B_2(u) &\;=\; B_1\!\left(\frac{49}{200}\right) + d\!\left(-\frac{23}{20} - \frac{133}{20} d + \frac{13}{10} d^2\right),
  \qquad d = u - \frac{49}{200},\quad u \in \left[\frac{49}{200}, \frac{51}{125}\right], \label{eq:B2}
\end{align}
extended constantly by $B_2(51/125)$ for $u > 51/125$. Direct evaluation gives $B_1(49/200) = B_2(49/200)$, $B_1'(49/200) \approx -0.8442 > -1.15 = B_2'(49/200)$ (strict convex corner), and $B_2(51/125) < 0$.

Along any smooth segment $r \notin \mathcal{J}$ where $v = B_i(u)$, \eqref{eq:tov-bondi} gives
\begin{equation}
  r(v - B_i(u))' \;=\; \frac{1}{1 - 2u}\Bigl[(1 - 2u)\bigl(2v - B_i'(u)(e - u)\bigr) - (e + v)(u + v)\Bigr],
  \label{eq:barrier-deriv}
\end{equation}
whose numerator is \emph{affine in $e \in [v, 3u]$}. Hence it is strictly negative for all $e \in [v, 3u]$ whenever both endpoint residuals are strictly negative:
\begin{align*}
  F_{\mathrm{core}}(B_i; u) &\;\equiv\; (1 - 2u)\bigl(2 B_i + (u - B_i) B_i'\bigr) - 2 B_i(u + B_i) \;<\; 0, \\
  F_{\mathrm{const}}(B_i; u) &\;\equiv\; (1 - 2u)\bigl(2 B_i - 2 u B_i'\bigr) - (3u + B_i)(u + B_i) \;<\; 0.
\end{align*}
Factoring out powers of $u$ on piece~1 gives $F_{\mathrm{core}}(B_1; u) = -u P_{1c}(u)$, $F_{\mathrm{const}}(B_1; u) = -u^2 P_{1k}(u)$, $F_{\mathrm{core}}(B_2; u) = -P_{2c}(u)$, $F_{\mathrm{const}}(B_2; u) = -P_{2k}(u)$, and $B_1(u) - 3u = u P_0(u)$. Exact rational Bernstein decompositions of the form
\[
  P(x) \;=\; \mu + \sum_{k=0}^{n} \gamma_k (x - a)^k (b - x)^{n-k}, \qquad \mu > 0,\;\; \gamma_k \ge 0
\]
certify $P_{1c} > 0$ (2 subintervals) and $P_{1k} > 0$ (1 interval) on $[0, 49/200]$, $P_{2c} > 0$ (1 interval) and $P_{2k} > 0$ (3 subintervals) on $[49/200, 51/125]$, and $P_0 > 0$ (1 interval) on $[0, 1/200]$.

Near $r = 0^+$, $u(r) \to 0$ implies $u(r) \in (0, 1/200]$ for small $r > 0$, whence $v(r) \le e(r) \le 3 u(r) < B_1(u(r)) = B(u(r))$. Suppose for contradiction that $v(r) - B(u(r))$ reaches $0$ on $[a, R]$, and let $x \in (a, R]$ be the infimum of such radii (first crossing). Even if $x \in \mathcal{J}$ is a density jump, the finiteness of $\mathcal{J}$ gives a punctured left neighborhood $(l, x)$ free of jumps on which $u(y)$ and $v(y)$ remain close to $u(x), v(x)$ by continuity, so the strict endpoint inequalities persist on $(l, x)$ (and at the corner $u(x) = 49/200$, $B(u) = \min(B_1(u), B_2(u))$ locally, so whichever piece is active has strictly negative derivative just to the left of $x$). Thus $v(y) - B(u(y))$ is strictly decreasing on $(l, x)$, forcing $v(y) > B(u(y))$ just left of $x$---contradicting that $x$ is the first crossing. Therefore $v(r) < B(u(r))$ for all $r \in (0, R]$. Since $v(r) \ge 0$ everywhere while $B(u) < 0$ for all $u \ge 51/125$, we conclude $u(r) < 51/125$ and $s(r)^2 > 23/125$ on $(0, R]$.
\end{proof}

\subsection{Part C: Proof of the Kerr Moment-of-Inertia Lower Bound}

\begin{theorem}[Equation-of-state-independent Kerr inertia lower bound]
\label{thm:kerr-inertia}
Every star in \texttt{MonotoneIsotropicDECStar} satisfies
\begin{equation}
  I \;=\; \frac{R^3 Z(R)}{6 s(R) + 2 Z(R)} \;>\; 4 M^3 \;=\; I_{\mathrm{Kerr}}.
  \label{eq:main-thm}
\end{equation}
In particular, by Lemma~\ref{lem:mean-density}, \eqref{eq:main-thm} holds for every piecewise-smooth regular isotropic DEC star (\texttt{PiecewiseRegularIsotropicDECStar}) whose energy density $\eps$ is non-increasing on $(0, R]$.
\end{theorem}

\begin{proof}
Set $\alpha = 69/100$ and $\Ibar_0 = 4$ in Proposition~\ref{prop:partA}, writing $W(r) \equiv W_{69/100}(r)$ and $D(r) \equiv W(r) - Z_c(4, s(r))$. The proof proceeds in four steps.

\medskip\noindent\textbf{Step 1 (Global non-negativity of the matter residual).}
At $\alpha = 69/100$ and $\Ibar_0 = 4$, direct polynomial algebra yields the exact factorizations in $t = s^2$:
\begin{align}
  N_e\!\left(\tfrac{69}{100}, 4, t\right) &\;=\; \frac{4}{25}\,t^2\,e_5(t), \qquad
  e_5(t) \;\equiv\; 31 t^5 - 117 t^4 + 51 t^3 - 1023 t^2 + 1737 t - 279, \label{eq:e5}\\
  N_e\!\left(\tfrac{69}{100}, 4, t\right) + N_p\!\left(\tfrac{69}{100}, 4, t\right)\!\bigl(4 - 4(1-t)^3\bigr) &\;=\; \frac{8}{25}\,t^3\,s_4(t), \qquad
  s_4(t) \;\equiv\; 31 t^4 - 48 t^3 - 363 t^2 + 498 t - 63. \label{eq:s4}
\end{align}
Each of $e_5(t)$ and $s_4(t)$ is certified strictly positive on the entire interval $t \in [23/125, 1]$ by a \emph{single-interval} exact rational Bernstein identity (no subdivision required):
\begin{align*}
  e_5(t) &\;=\; \frac{188093445758}{30517578125} + \sum_{k=1}^{5} \gamma_k^{(e)} \left(t - \tfrac{23}{125}\right)^k (1 - t)^{5-k} \;>\; 0, \\
  s_4(t) &\;=\; \frac{3925485571}{244140625} + \sum_{k=1}^{4} \gamma_k^{(s)} \left(t - \tfrac{23}{125}\right)^k (1 - t)^{4-k} \;>\; 0,
\end{align*}
with all rational coefficients $\gamma_k^{(e)}, \gamma_k^{(s)} > 0$ (explicitly listed in \texttt{e5\_pos\_0} and \texttt{s4\_pos\_0} of \texttt{T4A18\_KerrInertia.lean}). By Proposition~\ref{prop:partB}, $s(r)^2 \in (23/125, 1]$ on $(0, R]$, so $c_e > 0$ and $c_e + c_p > 0$ everywhere on $(0, R]$. Together with $c_{pp} = 4(69/100)(31/100)/s^3 = 2139/(2500 s^3) > 0$ and DEC ($0 \le \pi_p \le \pi_e$), the matter residual
\[
  \operatorname{Res}(r) \;\equiv\; c_e \pi_e + c_p \pi_p + c_{pp} \pi_p^2
\]
satisfies $\operatorname{Res}(r) \ge \min(c_e, c_e + c_p)\,\pi_e \ge 0$ everywhere on $(0, R]$, with strict inequality $\operatorname{Res}(r) > 0$ whenever $\eps(r) > 0$.

\medskip\noindent\textbf{Step 2 (Central non-negativity and strict interior entry).}
Since $\mathcal{J}$ is finite and $u, \pi_p, Z \to 0$ as $r \to 0^+$, there exists $a \in (0, R]$ such that $(0, a]$ contains no jumps and satisfies $u(r) \le 1/8$, $Z(r) \ge -1/8$ on $(0, a]$. On $(0, a]$, the linear coefficient \eqref{eq:L-def} obeys the elementary upper bound
\[
  L(r) \;=\; \frac{-Z s - Z_c s + u + (1 - 2\alpha)\pi_p}{s^2} - 3 \;\le\; \frac{s/8 + 1/8}{3/4} - 3 \;\le\; -2
\]
(since $1 - 2\alpha = -19/50 < 0$ and $s^2 \ge 3/4$). Consequently, if $D(r) < 0$ at any $r \in (0, a]$, then $r D'(r) = L D + \operatorname{Res} \ge -2 D > 0$. By the weighted central limits, $D(r) \to 0$ as $r \to 0^+$. A standard last-excursion argument shows $D(r) \ge 0$ for all $r \in (0, a]$: if $D(x) < 0$ for some $x \in (0, a]$, choosing $b \in (0, x)$ with $D(b) > D(x)$ and considering the supremum of points in $[b, x]$ where $D \ge (D(x) + \min(D(b), 0))/2$ produces a contradiction with $D' \ge 0$ whenever $D < 0$.

Moreover, $D$ cannot vanish identically on $(0, a]$: since $m(a) > 0$ and $m(0) = 0$, the mean value theorem yields some $x_1 \in (0, a)$ with $m'(x_1) = 4\pi x_1^2 \eps(x_1) > 0$, hence $\eps(x_1) > 0$ and $\operatorname{Res}(x_1) > 0$. If $D(x_1) = 0$, then $x_1 D'(x_1) = \operatorname{Res}(x_1) > 0$, so $D(x_0) > 0$ at some $x_0 \in (x_1, a)$. Thus there exists $x_0 \in (0, R]$ with $D(x_0) > 0$.

\medskip\noindent\textbf{Step 3 (Exponential barrier propagation to the surface across jumps).}
Since $L(r)/r$ is continuous on the compact interval $[x_0, R]$, it attains a minimum there; set $K \equiv 1 - \min_{[x_0, R]} (L(r)/r)$, so $K + L(r)/r \ge 1$ on $[x_0, R]$. Define $\eta \equiv D(x_0)/2 > 0$, $\varphi(r) \equiv e^{-K(r - x_0)}$, and $f(r) \equiv \eta \varphi(r) - D(r)$. Then $f(x_0) = -\eta < 0$ and $f$ is continuous on $[x_0, R]$. At any putative first zero $x \in (x_0, R]$ of $f$ (whether or not $x \in \mathcal{J}$), on a jump-free punctured left neighborhood of $x$ we have $D'(y) \ge (L(y)/y) D(y)$ (since $\operatorname{Res}(y) \ge 0$), hence
\[
  f'(y) \;\le\; -K \eta \varphi(y) - \frac{L(y)}{y}(\eta\varphi(y) - f(y))
  \;=\; -\Bigl(K + \frac{L(y)}{y}\Bigr)\eta\varphi(y) + \frac{L(y)}{y} f(y)
  \;\le\; -\eta\varphi(y) + \frac{L(y)}{y} f(y).
\]
Since the right-hand side is continuous and equals $-\eta\varphi(x) < 0$ at $x$, it stays strictly negative on a punctured left neighborhood of $x$, forcing $f(y) > 0$ just left of $x$ and contradicting that $x$ is the first crossing. Thus $f(R) < 0$, i.e.\ $D(R) > \eta \varphi(R) > 0$.

\medskip\noindent\textbf{Step 4 (Surface evaluation).}
Because $p(R) = 0$, we have $\pi_p(R) = 0$ and $Z(R) = W(R) > Z_c(4, s(R)) \ge 0$. Letting $u_R = M/R$ and $s_R = s(R)$, we have $1 - 8 u_R^3 > 0$ and $Z_c(4, s_R)(1 - 8 u_R^3) = 24 s_R u_R^3$. Multiplying $Z(R) > Z_c(4, s_R)$ by $1 - 8 u_R^3 > 0$ gives $Z(R) > 4 u_R^3 (6 s_R + 2 Z(R))$, and multiplying by $R^3/(6 s_R + 2 Z(R)) > 0$ yields $I = R^3 Z(R)/(6 s_R + 2 Z(R)) > 4 (u_R R)^3 = 4 M^3$.
\end{proof}

\section{Non-Vacuity: Explicit Analytic Smooth and Two-Layer Witnesses}
\label{sec:nonvacuity}

To demonstrate that neither \texttt{RegularIsotropicDECStar} nor \texttt{PiecewiseRegularIsotropicDECStar} is vacuous, we construct two explicit physical stars satisfying every hypothesis of Theorem~\ref{thm:kerr-inertia}.

\subsection{Smooth Witness: Homogeneous Schwarzschild Interior (\texorpdfstring{$C = 1/10$}{C = 1/10})}

Fix $R = 1$, $M = C = 1/10$, $\rho = 3C/(4\pi) = 3/(40\pi)$, $a = \sqrt{1 - 2C} = \sqrt{4/5}$, and $s(r) = \sqrt{1 - 2 C r^2}$. The exact Schwarzschild interior solution on $[0, 1]$ is
\begin{equation}
  m(r) \;=\; C r^3, \qquad \eps(r) \;=\; \rho, \qquad
  p(r) \;=\; \rho\,\frac{s(r) - a}{3a - s(r)}, \qquad
  N(r) \;=\; \frac{3a - s(r)}{2}.
  \label{eq:schw-interior}
\end{equation}
Since $a = \sqrt{4/5} > 1/2$, we have $3a - s(r) \ge 3a - 1 > 0$, $s(r) \ge a > 0$, $p(1) = 0$, and $0 \le p(r)/\rho \le (1 - a)/(3a - 1) \approx 0.0627 < 1$, verifying DEC and $4\pi r^3 \eps = 3 m$ identically. Direct differentiation checks $m' = 4\pi r^2 \rho$ and $p' = -4\rho a C r / [s(3a-s)^2]$, matching \eqref{eq:tov}.

For Hartle's equation \eqref{eq:hartle-2nd}, $j(r) = s(r)/N(r) = 2 s(r)/(3a - s(r))$ is positive, even, and real-analytic on $[0, 1]$, with $j'(r) = 6 a s'(r)/(3a-s)^2 < 0$ for $r \in (0, 1]$. The regular singular point $r=0$ of \eqref{eq:hartle-2nd} has indicial roots $0$ and $-3$, so Frobenius theory yields a unique smooth solution $\wbar$ normalized by $\wbar(0) = 1$, $\wbar'(0) = 0$, with $\wbar(r) = 1 + \frac{8\pi}{5}(\rho + p(0)) r^2 + O(r^4)$ \cite{Petroff2007}. Integrating \eqref{eq:hartle-2nd} from $0$ to $r$ gives
\[
  r^4 j(r)\,\wbar'(r) \;=\; \int_0^r (-4 t^3 j'(t))\,\wbar(t)\,\dd t.
\]
While $\wbar > 0$, the integrand is strictly positive on $(0, r]$, so $\wbar' > 0$ and $\wbar$ is strictly increasing from $1$; hence $\wbar$ can never vanish on $[0, 1]$. Therefore $Z(r) = r s(r)\wbar'(r)/\wbar(r)$ is smooth on $[0, 1]$, solves \eqref{eq:hartle-Z} globally, satisfies $Z(0) = 0$ and $Z(r) = O(r^2)$ as $r \to 0^+$, and furnishes a complete smooth witness. Numerical integration gives $Z(1) = 0.25824624$ and $\Ibar = I/M^3 = 43.896617 > 4$.

\subsection{Discontinuous Witness: Two-Layer DEC Star (\texorpdfstring{$\mathcal{J} = \{1/2, 1\}$}{J = \{1/2, 1\}})}

Set $R = 1$, $\mathcal{J} = \{1/2, 1\}$, and $\rho = 3/(800\pi)$, with piecewise-constant non-increasing density
\[
  \eps(r) \;=\; \begin{cases}
    2\rho, & 0 < r < \tfrac{1}{2},\\
    \rho, & \tfrac{1}{2} \le r < 1,\\
    0, & r \ge 1.
  \end{cases}
\]
Integrating $m' = 4\pi r^2 \eps$ with $m(0) = 0$ gives a continuous mass profile satisfying $0 < m(r) \le r^3/100$ on $(0, 1]$, so $s(r) \ge \sqrt{49/50} > 0$ and $m(r)/r \to 0$. Integrating the TOV equation \eqref{eq:tov} inward from $p(1) = 0$ with continuous pressure across $r = 1/2$, while $0 \le p \le \rho$ we have $m + 4\pi r^3 p \le r^3/100 + 3 r^3/200 = r^3/40$, $s^2 \ge 49/50$, and $\eps + p \le 3\rho$, hence
\[
  0 \;\le\; -p'(r) \;\le\; \frac{3\rho\,(r^3/40)}{r^2\,(49/50)} \;=\; \frac{15}{196}\,\rho r
\]
on each open layer. Integrating over $[0, 1]$ yields $0 \le p(r) \le (15/392)\rho < \rho \le \eps(r)$ on $(0, 1)$, so DEC holds strictly on both layers and at the interface $r = 1/2$ (and $p(1) = 0 = \eps(1)$ at the surface). Matching the regular core Hartle solution ($\wbar(0)=1$) continuously in $\wbar$ and $F \equiv r^4 j \wbar'$ at $r = 1/2$ yields $F' = -4 r^3 j' \wbar > 0$ on both open layers, so $\wbar \ge 1$ and $Z = s F/(r^3 j \wbar)$ is continuous on $(0, 1]$ with $Z(r) = O(r^2)$ at the center. This verifies every field of \texttt{PiecewiseRegularIsotropicDECStar} together with \texttt{AntitoneOn S.$\eps$ (Ioc 0 S.R)}.

\begin{remark}
While the proofs above establish non-vacuity mathematically, formalizing \texttt{Nonempty MonotoneIsotropicDECStar} inside Lean~4 requires an internal Mathlib existence proof for the regular Frobenius solution of Hartle's ODE \eqref{eq:hartle-2nd}. That formalization is a well-defined next step and is not yet included in the repository.
\end{remark}

\section{Sharpness, Beyond-DEC Extension, and Necessity of Isotropy}
\label{sec:sharpness}

\subsection{Numerical Sharpness Benchmarks}

Theorem~\ref{thm:kerr-inertia} proves the strict Kerr benchmark $\Ibar = I/M^3 > 4$, but $4$ is \emph{not} the sharp infimum of $\Ibar$ over \texttt{MonotoneIsotropicDECStar}. Table~\ref{tab:sharpness} compares high-precision numerical integrations of the coupled Hartle--TOV equations across two extremal families:
\begin{enumerate}
  \item \emph{Uniform-density Schwarzschild stars} ($\eps = \text{const}$): reach the DEC ceiling $p_c = \eps$ at $C = 3/8 = 0.375$ with $\Ibar = 4.654887$, and approach the Buchdahl limit $C \to 4/9$ (violating DEC with $p_c/\eps \to \infty$) with $I/(MR^2) \to 0.7992$ and $\Ibar \to 4.0459$ \cite{ChandrasekharMiller1974,Posada2017}.
  \item \emph{Two-zone monotone isotropic DEC stars} ($p = \eps$ stiff core on $[0, r_{\mathrm{sw}}]$ matched continuously in $\eps$ and $p$ to a constant-density envelope $\eps(r) = \eps(r_{\mathrm{sw}})$ on $[r_{\mathrm{sw}}, R]$, in units $\eps_c = 1$): satisfy $0 \le p \le \eps$ and $\eps' \le 0$ everywhere, exceed the uniform DEC compactness limit $3/8$ to reach $C = 0.398800$ ($2C = 0.797600$), and attain a minimum $\Ibar = 4.285459$ near $r_{\mathrm{sw}} = 0.40$.
\end{enumerate}

\begin{table}[ht]
\centering
\small
\begin{tabular}{@{}lrrrr@{}}
\toprule
Configuration & $p_c/\eps_c$ & $C = M/R$ & $I/(MR^2)$ & $\Ibar = I/M^3$ \\
\midrule
Uniform Schwarzschild ($C = 0.10$) & $0.06272$ & $0.100000$ & $0.438966$ & $43.896617$ \\
Uniform Schwarzschild ($C = 0.30$) & $0.40958$ & $0.300000$ & $0.566436$ & $6.293731$ \\
Uniform Schwarzschild (DEC limit $C = 3/8$) & $1.00000$ & $0.375000$ & $0.654593$ & $4.654887$ \\
Two-zone DEC ($r_{\mathrm{sw}} = 0.25$) & $1.00000$ & $0.398128$ & $0.686416$ & $4.330515$ \\
Two-zone DEC ($r_{\mathrm{sw}} = 0.30$, near max $C$) & $1.00000$ & $0.398800$ & $0.684202$ & $4.302025$ \\
Two-zone DEC ($r_{\mathrm{sw}} = 0.40$, near min $\Ibar$) & $1.00000$ & $0.395793$ & $0.671326$ & $4.285459$ \\
\midrule
Uniform Schwarzschild ($C = 0.4444$, violates DEC) & $1666.7$ & $0.444400$ & $0.799034$ & $4.045920$ \\
Kerr black hole (slow-rotation limit) & --- & $0.500000$ & --- & $4.000000$ \\
\bottomrule
\end{tabular}
\caption{Coupled Hartle--TOV benchmarks for uniform-density and two-zone monotone isotropic DEC stars (all computed with \texttt{scipy.integrate.solve\_ivp}, \texttt{rtol}=$10^{-12}$--$10^{-13}$). Two-zone monotone DEC stars exceed $C = 3/8$ while remaining well above $\Ibar = 4$.}
\label{tab:sharpness}
\end{table}

Determining the exact sharp infimum $\inf \Ibar$ over \texttt{MonotoneIsotropicDECStar} (which Table~\ref{tab:sharpness} places in $[4, 4.2855]$) remains an interesting variational problem.

\subsection{Beyond-DEC Extension and the Open Buchdahl Question}

Optimizing the shift parameter $\alpha$ in Proposition~\ref{prop:partA} at $\Ibar_0 = 4$ yields $\alpha_* \approx 0.69664743$, which keeps both $c_e > 0$ and $c_e + c_p > 0$ up to $u \le 0.41149742$. Combining this with Bondi's relaxed compactness ceiling $\mu_{\max}(w)$ for $0 \le p \le w \eps$ extends the barrier proof of $I > 4 M^3$ beyond DEC up to $w_* \approx 1.155605$ (where $\mu_{\max}(w_*) \approx 0.411339$).

By contrast, dropping the upper pressure bound $p \le w\eps$ altogether (the full Buchdahl class: $p \ge 0$, $\eps' \le 0$, $p_c \to \infty$) defeats the linear shift \eqref{eq:W-D-def}: for $u \in (0.3674, 4/9)$, $c_p(\alpha, 4, s)$ is negative whenever $c_e(\alpha, 4, s) \ge 0$, so arbitrarily large $p/\eps$ can make $c_e \pi_e + c_p \pi_p$ negative. Whether $I > 4 M^3$ holds across the entire Buchdahl class (as suggested by the uniform-density Buchdahl limit $\Ibar \to 4.046$) therefore requires a different Lyapunov functional and remains open.

\subsection{Necessity of the Isotropy Hypothesis}

If pressure isotropy ($p_r = p_t$) is relaxed to allow anisotropic tangential stress $p_t \neq p_r$ while preserving the full covariant dominant energy condition ($|p_r| \le \eps$, $|p_t| \le \eps$, $\eps \ge 0$), the inequality $I > 4 M^3$ \emph{fails}. Indeed, the anisotropic generalization of Proposition~\ref{prop:partA} introduces the tangential stress term $+4(1-\alpha)\pi_{p_t}/s$ into the Hartle residual; strong tangential tension ($p_t < 0$) reduces the frame-dragging source $4(\pi_e + \pi_{p_t})/s$ and lowers $I/M^3$. Explicit thin-shell and smooth weak-gravity anisotropic DEC configurations achieve $\Ibar \in [3.74, 3.77] < 4$ (see \cite{BartaGergely2026} for isotropic EOS benchmarks and the companion script \texttt{scripts/verify\_t4a18\_kerr\_inertia\_floor.py} for anisotropic diagnostics). Consequently, the isotropy assumption in Theorem~\ref{thm:kerr-inertia} cannot be dropped without replacing it by a lower bound on $p_t$.

\section{Lean~4 Formalization, Verification, and Workflow Disclosure}
\label{sec:lean}

The complete proof of Lemma~\ref{lem:mean-density}, Proposition~\ref{prop:partA}, Proposition~\ref{prop:partB}, and Theorem~\ref{thm:kerr-inertia} is formalized in Lean~4.36.0 + Mathlib in \texttt{NaturalMechanics/T4A18\_KerrInertia.lean} ($1{,}576$ lines) in the public repository \url{https://github.com/iancenco/relativistic-bounds-lean4}. Key verification properties are:
\begin{itemize}
  \item \emph{Zero floating-point arithmetic:} All polynomial positivity lemmas (\texttt{P1c\_pos}, \texttt{P1k\_pos}, \texttt{P2c\_pos}, \texttt{P2k\_pos}, \texttt{P0\_pos}, \texttt{e5\_pos}, \texttt{s4\_pos}) are verified as exact rational identities by \texttt{ring} and \texttt{positivity}.
  \item \emph{Standard axioms only:} Continuous integration (\texttt{.github/workflows/lean.yml}) runs \texttt{lake exe cache get}, \texttt{lake build}, a prohibited-token scan (\texttt{sorry}, \texttt{admit}, custom \texttt{axiom}, \texttt{unsafe}, \texttt{implemented\_by}, \texttt{native\_decide}, \texttt{@[extern]}), and an automated \texttt{\#print axioms} check verifying that every headline theorem depends solely on \texttt{[propext, Classical.choice, Quot.sound]}.
  \item \emph{Main Lean~4 declarations:}
    \begin{itemize}
      \item \texttt{MonotoneIsotropicDECStar.t4a18\_unconditional\_kerr\_inertia\_floor} (Theorem~\ref{thm:kerr-inertia}, mean-density form);
      \item \texttt{t4a18\_kerr\_inertia\_floor\_of\_antitone\_piecewise} (Theorem~\ref{thm:kerr-inertia}, literal non-increasing $\eps$ with finitely many jumps);
      \item \texttt{MonotoneIsotropicDECStar.compactness\_wall} (Proposition~\ref{prop:partB}, auxiliary compactness bound $m/r < 51/125$);
      \item \texttt{partA} (Proposition~\ref{prop:partA}, exact shifted Hartle--TOV Riccati identity).
    \end{itemize}
\end{itemize}

\paragraph{AI-assisted workflow disclosure.}
The continuum barrier search, exact rational Bernstein certificate generation, Lean~4 formalization, and independent statement-level audit (\texttt{AUDIT-T4A18.md}) were carried out using an AI-assisted multi-model workflow directed by the author. In particular, an adversarial statement audit caught an initial omission of interior density jumps in the Lean star structure, prompting the jump-inclusive \texttt{PiecewiseRegularIsotropicDECStar} revision prior to publication of the repository. Logical correctness of the formalized theorems depends strictly on the Lean~4 kernel and standard axioms; physical faithfulness of the encoded assumptions and novelty relative to the literature should be evaluated directly from this manuscript and the accompanying source files.

\begin{thebibliography}{99}

\bibitem{Hartle1967}
J.~B.~Hartle,
``Slowly rotating relativistic stars. I. Equations of structure,''
\emph{Astrophys.\ J.} \textbf{150} (1967) 1005--1029.

\bibitem{Buchdahl1959}
H.~A.~Buchdahl,
``General relativistic fluid spheres,''
\emph{Phys.\ Rev.} \textbf{116} (1959) 1027--1034.

\bibitem{Bondi1964}
H.~Bondi,
``Massive spheres in general relativity,''
\emph{Proc.\ R.\ Soc.\ Lond.\ A} \textbf{282} (1964) 303--317.

\bibitem{ChandrasekharMiller1974}
S.~Chandrasekhar and J.~C.~Miller,
``On slowly rotating homogeneous masses in general relativity,''
\emph{Mon.\ Not.\ R.\ Astron.\ Soc.} \textbf{167} (1974) 63--80.

\bibitem{SabbadiniHartle1977}
A.~G.~Sabbadini and J.~B.~Hartle,
``Upper and lower bounds on the moments of inertia of relativistic, spherical, perfect fluid neutron stars,''
\emph{Ann.\ Phys.} \textbf{104} (1977) 95--134.

\bibitem{Hartle1978}
J.~B.~Hartle,
``Bounds on the mass and moment of inertia of non-rotating neutron stars,''
\emph{Phys.\ Rep.} \textbf{46} (1978) 201--247.

\bibitem{BarracoHamity2002}
D.~Barraco and V.~H.~Hamity,
``Maximum mass of a spherically symmetric isotropic star,''
\emph{Phys.\ Rev.\ D} \textbf{65} (2002) 124028.

\bibitem{Petroff2007}
D.~Petroff,
``The homogeneous, slowly rotating fluid sphere in general relativity,''
\emph{Class.\ Quantum Grav.} \textbf{24} (2007) 1055--1066 [arXiv:gr-qc/0701081].

\bibitem{KarageorgisStalker2008}
P.~Karageorgis and J.~G.~Stalker,
``Sharp bounds on $2m/r$ for static spherical objects,''
\emph{Class.\ Quantum Grav.} \textbf{25} (2008) 195021 [arXiv:0707.3632].

\bibitem{Andreasson2008}
H.~Andr\'easson,
``Sharp bounds on $2m/r$ of general spherically symmetric static objects,''
\emph{J.\ Differential Equations} \textbf{245} (2008) 2243--2266.

\bibitem{YagiYunes2013}
K.~Yagi and N.~Yunes,
``I-Love-Q relations in neutron stars and their applications to astrophysics, gravitational waves, and fundamental physics,''
\emph{Phys.\ Rev.\ D} \textbf{88} (2013) 023009.

\bibitem{BreuRezzolla2016}
C.~Breu and L.~Rezzolla,
``Maximum mass, moment of inertia and compactness of relativistic stars,''
\emph{Mon.\ Not.\ R.\ Astron.\ Soc.} \textbf{459} (2016) 646--656.

\bibitem{Posada2017}
C.~Posada,
``Slowly rotating supercompact Schwarzschild stars,''
\emph{Mon.\ Not.\ R.\ Astron.\ Soc.} \textbf{468} (2017) 2128--2139.

\bibitem{BartaGergely2026}
D.~Barta and L.~\'A.~Gergely,
``Compactness relations and Hartle--Thorne quadrupoles for neutron-star spin flip-flops,''
\emph{Universe} (November 2026).

\end{thebibliography}

\end{document}
